\section{Mission scoring}
\label{sec:mission}

A task tree reports events on \code{/task/events}. Its rubric gives required
counts $n_e$, event points $p_e$ and count caps $c_e$ (default one). For
observed counts $k_e$ the event score is
\[
 S_{\rm events}=\sum_e p_e\min(k_e,c_e),\qquad
 P_{\rm events}=\sum_e p_ec_e.
\]
The last \code{task\_finished} supplies the outcome, falling back to the process
outcome or INCOMPLETE. Elapsed time comes from that finish event, otherwise the
last event, defaulting to zero. PASS requires SUCCESS, $k_e\ge n_e$ for every
required event, and elapsed time $\le T$ when a time limit is specified.
The procedure stores this verdict and the task referee verdict separately.
\file{mhseals\_testing/analysis/mission.py; procedures/mission.py; config/missions/}

\subsection{The task referee}

The task referee uses task-tagged course truth and the boat's true track;
absent truth or track yields an explicit N/A. Pair and circle tests also yield
N/A when truth contains no elements tagged for that test.

\paragraph{Gates.} Truth markers are paired greedily in increasing distance,
each used once, with $0.5<w\le w_{\max}$ metres. The maximum is $8$\,m for
channel and pair tests and $4$\,m otherwise; return-home gates pair black
markers with each other. For a gate segment $(\bm a,\bm b)$ and track step
$(\bm p,\bm q)$ define
\[
 c(\bm o,\bm u,\bm v)=(u_x-o_x)(v_y-o_y)-(u_y-o_y)(v_x-o_x),
 \quad(d_1,d_2,d_3,d_4)=
 (c(\bm a,\bm b,\bm p),c(\bm a,\bm b,\bm q),
 c(\bm p,\bm q,\bm a),c(\bm p,\bm q,\bm b)).
\]
A crossing occurs exactly when
\[
 d_1\ne0,\qquad d_1d_2\le0,\qquad d_3d_4\le0;
 \qquad \text{direction}=-\operatorname{sign}(d_1).
\]
Reaching the gate line or an endpoint counts on the arriving step, at the
second sample's time; a step starting on the line does not count. Pair tests
require at least two crossings of one gate with different first and last
directions. Channel, follow-path and return-home tasks require a nonempty
set of gates and at least one crossing of every gate. Follow-path also requires
the last reported species count to equal the number of yellow truth buoys.

\paragraph{Contact.} Rotate an object's offset into hull coordinates $(a,c)$.
For length $L=1.063$\,m, beam $B=0.895$\,m and object radius $r_b$
(default $0.2$\,m), clearance is
\[
 \delta=\sqrt{\max(|a|-L/2,0)^2+\max(|c|-B/2,0)^2}-r_b.
\]
Only sampled poses with $\delta<0$ count as contact. Every task requires no
contact with its labelled objects; unlabelled dock structure is excluded.

\paragraph{Circling.} For bearings $\theta_k$ about a mark, the swept angle is
\[
 A=\frac{180}{\pi}\left|\sum_{k\in\mathcal I}
       \operatorname{remainder}(\theta_k-\theta_{k-1},2\pi)\right|.
\]
$\mathcal I$ includes only consecutive finite samples within $8$\,m of the
mark. Leaving this radius breaks angle continuity but preserves the accumulated
sum. A circle test requires $A\ge345^\circ$ about some yellow buoy and
$\|\bm p_{\rm last}-\bm p_{\rm first}\|\le0.75$\,m. A speed challenge uses
the first paired gate and first blue mark: first and last gate-crossing
directions must differ, the rounded swept angle must be at least $300^\circ$,
and the last spill count must equal the number of black truth buoys.
Its gate-to-gate time is the last crossing time minus the first.

\paragraph{Docking.} In a bay's frame, with centre offset $(a,c)$ and opening
yaw $\psi_b$, \code{held\_in\_bay} accumulates the longest uninterrupted
sampled hold satisfying
\[
 |a|\le D/2+0.3,\qquad |c|\le W/2,\qquad
 \cos(\psi-\psi_b-\pi)>\cos30^\circ.
\]
An invalid sample resets the hold. The matching-banner, unoccupied bays must
include a hold of at least \code{dock\_time} (default $5$\,s); every other bay
must have a hold shorter than $1$\,s.

\paragraph{Mission stages.} Each stage uses the track and events between its
first start and first completion or skip, inclusively; an absent terminal event
leaves the end open. Missing required starts fail. An optional unattempted or
skipped stage is SKIPPED unless its stage-window track contacts a labelled
object. The mission passes exactly when every stage is PASS or SKIPPED:
\[
 \text{PASS}_{\rm mission}\iff
 \bigwedge_s\bigl(\text{verdict}_s\in\{\text{PASS},\text{SKIPPED}\}\bigr).
\]
\file{mhseals\_testing/analysis/referee.py, geometry.py; procedures/mission.py}

\subsection{Ground-truth run points}

The run scorer uses course elements, ordered gates, bay geometry, state changes,
authority records and the true track. Task events do not establish earned
behaviour. It keeps earned points, provisional points and possible points
separately; report-only rows are unscorable when RoboCommand capture is absent.
For row points $s_i$, provisional amounts $q_i$ and $N_{\rm strike}$ strikes,
\[
 S=\sum_i s_i-25\min(N_{\rm strike},20),\qquad Q=\sum_i q_i,\qquad
 P=\sum_{i:\,\text{status}_i\ne\text{n/a}}P_i.
\]
Scorable possible points also exclude unscorable rows. The total is not floored
at zero. Missing behaviour is missed; behaviour needing an uncaptured report
is provisional rather than earned.

\paragraph{Window and gate passes.} The scorer clips finite track samples to
the recorded window, ending at the first intervention by \code{rc},
\code{keyboard}, \code{hold} or \code{external}, or an earlier sampled boundary
exit. When all headings are zero or absent it derives yaw from track motion.
Let $s_k=(\bm p_k-\bm m)\cdot\bm n$ for gate normal $\bm n$.
A pass requires a change between sides $s<-0.3$ and $s>0.3$\,m. The last
zero crossing is interpolated, including yaw by the wrapped angular difference:
\[
 f=\frac{s_{k-1}}{s_{k-1}-s_k},\qquad
 (t_z,\bm p_z)=(1-f)(t_{k-1},\bm p_{k-1})+f(t_k,\bm p_k).
\]
The centre must cross between the marker centres. With tangent
$\bm u=(-n_y,n_x)$, the hull's projected half-width is
\[
 h_u=\tfrac L2|\bm e_\psi\cdot\bm u|
       +\tfrac B2|\bm e_\psi^\perp\cdot\bm u|.
\]
Whole-hull fit requires the interval
$[\bm u\cdot(\bm p_z-\bm m)-h_u,\bm u\cdot(\bm p_z-\bm m)+h_u]$
to lie between the two buoys' inner edges. Ordered gate sequences require
strictly increasing crossing times and whole-hull fit at every gate.

Entry earns $100$ for a complete outbound entry-gate sequence without buoy
strikes in the window from $3$\,s before its first pass to $3$\,s after its last.
The first clean attempt sets $t_{\rm entry}$. The first subsequent complete
reverse sequence earns $200$ for exit and cuts off all later evidence.
Other run rows require $t_{\rm entry}<t\le t_{\rm cut}$; task-only scoring
omits the entry prerequisite. Entry and exit reports are $100$ possible each,
unscorable without report capture. Initial standoff $\ge1.8288$\,m before the
first gate is recorded, and a course-start offset above $10$\,m is a warning.

\paragraph{Debris clearance.} Ordered outbound and reverse channel sequences
earn $100$ each. Red avoidance is $100$ provisional when every designated red
is approached within $20$\,m, has strictly positive sampled minimum clearance,
and is not fully circled. Green rescue is $200$ provisional for a full
$360^\circ$ circle about the configured number of greens (default one).
The angular metric differs from the task test: within each uninterrupted stay
inside $8$\,m, form cumulative wrapped angles $z_k$ and retain the largest
absolute signed subinterval sweep,
\[
 A_{\rm run}=\frac{180}{\pi}\max_{\text{one stay},\,i\le j}|z_j-z_i|.
\]
Nonfinite or out-of-radius positions reset the stay. Reporting all debris has
$200$ possible, unscorable without report capture.

\paragraph{Speed challenge.} The first valid loop starts with a whole-hull
outbound gate pass and ends at the next inbound pass, also with whole-hull fit.
The intervening sampled path, closed by a straight segment, must have nonzero
integer winding about the turn mark:
\[
 w=\operatorname{round}\!\left[\frac1{2\pi}
 \sum_{\text{closed path}}\operatorname{remainder}(\theta_{k+1}-\theta_k,2\pi)\right].
\]
The loop earns $200$. Direction earns $150$ when the indicator at entry calls
for its sign: red means counter-clockwise, green clockwise. For gate-to-gate
time $T_g$, time points are
\[
 s_{\rm time}=\operatorname{round}(\max(0,250-T_g),1).
\]
The colour/time report has $100$ possible, unscorable without capture.

\paragraph{Supply, marina and harbor.} Each delivery type has possible points
$\sum_{j=1}^{\min(n,3)}(75,125,200)_j$; without delivery sensing these rows are
unscorable. The delivery report has $100$ possible.
Marina docking uses centre containment $|a|\le D/2$, $|c|\le W/2$, with
free heading. A hold reaches $2$\,s when successive inside samples are no more
than $0.5$\,s apart. The first held slip in the allowed window decides the tier:
occupied earns zero, unoccupied $50$, green $200$, or lowest-numbered unoccupied
green $350$. Only the highest tier applies; ties for lowest count, and a missing
sign ranks as $9$. The dock report has $100$ possible.

After the first recorded horn, the first outside-to-inside zone entry within
$120$\,s earns $150$, or $300$ for the commanded zone (one horn: zone 1;
two: zone 2). Zone radius defaults to $3$\,m. Presence at the horn is not an
entry, and task work on the way does not cancel credit. The confirmation report
has $150$ possible. A complete silent horn log means no alert; a planned horn
or partial recorded horn can yield only provisional zone points. Unknown or
contradictory horn evidence makes the zone row unscorable when it cannot decide
an alert.

\paragraph{Strikes and evidence.} The clearance formula above also detects
strike episodes, using the course's hull dimensions when supplied. An episode
starts at $\delta<0$ and rearms only after $\delta>0.2$\,m; dock spine, tine,
sign and indicator roles are excluded. Judge-reported contacts replace geometric
strikes when that capability is selected; an absent judge sheet makes strikes
unscorable. Indicator state starts from the truth program and recorded changes
override it. Partial or contradictory state logs make earned direction and dock
rows provisional. Surveyed truth is rescored after growing objects by
$\sigma$, shrinking slips by $2\sigma$ and zone radii by $\sigma$, and under
the opposite changes; a row whose result changes becomes provisional at the
more generous points. With no recorded time limit all earned rows become
provisional. The $900$\,s timeslot comparison is a warning, not a substitute
scoring window.

\paragraph{Run outcome.} A record carrying an outcome uses only its own file
links, including partial links after an abort or error; missing links never
fall back to the session's latest track or task-event files. A started run with
a linked track is scored from the available evidence. Explicit \code{started}
decides whether launch occurred; when absent, a journal start or a recorded
run-start timestamp establishes it. A journal stop takes precedence over the
record's \code{outcome}; without a stop the record supplies the terminal outcome.
For a started run, an outcome other than \code{complete} is a mission failure.
A run with no track and a terminal failure still counts as a failure; a run
that did not start is excluded from the failure rate. A completed run without
scorable evidence remains unknown. Failure flags cover missing exit, hold,
takeover, contact, false claim and mission failure. Any observed failure makes
a run unclean; absent failures with unknown flags leave cleanliness unknown. A run is known when all six
flags are known or it has a terminal harness failure; the failure rate is
\[
 \widehat p_{\rm fail}=\frac{\#\{\text{known runs with a failure}\}}
                              {\#\{\text{known runs}\}}.
\]
Unknown observations do not establish a clean run.
\file{mhseals\_testing/analysis/score.py, run\_record.py, reliability.py, geometry.py}
